\documentclass[10pt,letterpaper]{article}
\usepackage[margin=1in]{geometry}
\usepackage{iftex}
% mathptmx has no TU-encoded faces, so under XeTeX (Tectonic) version 1 fell back
% to Latin Modern regular and lost every bold and italic. Use Termes there.
\ifPDFTeX\usepackage{mathptmx}\else\usepackage{fontspec}\setmainfont{texgyretermes}[Extension=.otf,UprightFont=*-regular,BoldFont=*-bold,ItalicFont=*-italic,BoldItalicFont=*-bolditalic]\fi
\usepackage{amsmath,booktabs,natbib,microtype,parskip,enumitem}
\usepackage[colorlinks=true,linkcolor=blue,citecolor=blue,urlcolor=blue]{hyperref}
\setlength{\parindent}{0pt}
\setlength{\emergencystretch}{3em}

\title{\textbf{Weight Custody Manifest: Portable, Attestation-Gated Custody for Proprietary Model Weights}}
\author{
  Imran Siddique \\
  OPAQUE Systems \\[4pt]
  \normalsize \url{https://github.com/agentrust-io/weight-custody-manifest}
}
\date{October 1, 2026 \\ \small Technical report, version 2}

\begin{document}
\maketitle

\begin{abstract}
Deploying proprietary model weights into infrastructure controlled by a customer reverses the usual confidential-computing trust problem: the model builder, rather than the infrastructure owner, supplies the secret. Existing confidential-computing stacks can attest workloads and release secrets, while model-signing systems authenticate model artifacts. Neither capability alone specifies continuing custody of a particular weight artifact after release. This technical report describes the Weight Custody Manifest (WCM), a vendor-neutral protocol and conformance model that binds an exact weight digest, an approved workload measurement, attestation policy, a fresh transport key, renewable authorization, terminal refusal and wipe evidence, and derivative lineage. WCM distinguishes cryptographic custody against software adversaries from accountability-grade controls when the infrastructure operator physically owns the machine. The evaluated Python reference implementation, version 0.28.0, supplies 91 portable conformance vectors across four levels. Its test suite produced 618 passing tests and three skips on September 3, 2026, and the same counts when rerun at the same commit for this edition. Later findings bound those results. The evaluated key broker compared the approved workload measurement with an unsigned evidence field rather than the signed report, and did not require GPU confidential-compute mode before release; both are fixed in later releases. The SEV-SNP platform used for the hardware run does not enable the ciphertext hiding that the cryptographic-custody claim against a hypervisor-privileged operator requires. The contribution is a portable artifact-to-runtime custody contract with explicit lifecycle evidence, derivative accountability, and conformance, not a new attestation primitive or a certification of any deployment.
\end{abstract}

\paragraph{Scope of this edition.}
This is an author-prepared technical report, not peer reviewed and not a
normative specification. It revises a manuscript dated September 3, 2026. The
software results were rerun on September 27, 2026 at the evaluated commit
(\texttt{2acedfa}, WCM 0.28.0) with current dependency resolutions; the hardware
runs were not repeated. Findings published after the evaluation are reported in
Section~\ref{sec:after} beside the results they qualify. Current behavior,
normative requirements, and limitations are maintained in the project repository.
Version 2 (October 1, 2026) attributes the attestation break on patched SEV-SNP
systems to Battering RAM rather than BadRAM, adds the verification defects fixed
in WCM 0.28.5 and 0.29.0 to Section~\ref{sec:after}, states that sealed release
is a broker setting, restores bold and italic type lost in the version 1 build,
and corrects bibliography metadata. The software results were rerun again on
October 1, 2026 with the same counts.

\noindent\textbf{Patent Pending.}

\section{Introduction}

Frontier-model deployment increasingly separates the model builder from the infrastructure operator. A builder may need to run a closed model inside a customer's cloud account, sovereign environment, or on-premises cluster without giving the operator a reusable plaintext copy of the weights. Encryption at rest is insufficient because an inference runtime eventually needs a decryption capability. Remote attestation and key brokers address the initial release decision: a verifier appraises evidence from a protected workload, and a relying party releases a secret only when policy is satisfied \citep{rats,trustee,nvidia2026}.

The initial release is only one boundary in a longer custody lifecycle. A deployer must identify the exact model bytes, bind them to the intended workload and policy, prevent replay of old evidence, communicate the released key only to the attested instance, renew or revoke authority, stop serving when authority lapses, preserve tamper-evident evidence, and account for derived artifacts. Existing systems implement important subsets, but their deployment-specific policy and secret-delivery interfaces do not constitute a portable model-custody record.

WCM contributes a composition rather than a new attestation primitive. Its signed manifest is the stable object that a model builder and independent verifier can share across providers. It names exact artifacts and allowed measurements, constrains key release, and connects the release decision to runtime custody and lineage records. Four design rules follow:

\begin{enumerate}[leftmargin=*]
  \item an artifact signature is necessary but does not authorize execution;
  \item attestation evidence is necessary but does not identify which weight artifact may be released;
  \item a successful release is temporary, so continued use requires renewable authority and observable terminal behavior;
  \item physical ownership of current confidential-computing hardware is outside the strongest cryptographic guarantee.
\end{enumerate}

\section{Threat model and guarantee boundary}

The protected asset is the plaintext weight artifact and its release key. The model builder controls the artifact digest, signing key, acceptable workload measurements, release policy, and revocation decisions. The customer or platform operator controls scheduling, storage, networking, host software, and availability. The protected runtime and verifier form part of the trusted computing base; the host operating system and remote software administrator need not.

WCM defines two non-interchangeable claims. \emph{Cryptographic custody} targets remote and software adversaries, including a compromised host stack, provided the selected TEE and its verified configuration actually defend that adversary. \emph{Accountability-grade custody} is the maximum claim against an operator with physical possession of the machine. TEE.fail demonstrates inexpensive DDR5 interposition attacks against contemporary TDX and SEV-SNP systems and shows consequences for GPU confidential computing \citep{teefail}. BadRAM broke SEV-SNP attestation through memory aliasing \citep{badram}, and Battering RAM restores that break on systems carrying the firmware fix by aliasing memory at run time with an interposer \citep{batteringram}. WCM therefore records the physical-control tier and residual risk rather than transforming a signed quote into a claim of physical impossibility.

Denial of service, malicious model code, side channels outside the selected profile, compromised builder keys, and truthful-but-misleading application outputs remain out of scope. WCM does not certify an entire organization's RAND security level; RAND levels describe whole-system postures, not one control \citep{randweights}.

\section{Protocol}

\subsection{Manifest and artifact identity}

The model builder computes a deterministic digest over the weight artifact and signs a manifest containing the artifact identity, allowed workload measurements, evidence requirements, release-key identifier, validity interval, renewal policy, and lineage policy. Model signing systems such as OpenSSF Model Signing establish artifact authenticity and integrity \citep{oms}; WCM consumes that result and adds a runtime authorization decision. A valid artifact signature does not imply that the current workload may decrypt the artifact.

\subsection{Fresh attestation and channel binding}

The key broker issues a fresh single-use challenge. The protected workload generates an ephemeral transport key and obtains CPU, and when required GPU, evidence that binds the challenge, workload state, and transport key. The verifier validates certificate chains, freshness, revocation, measurements, and policy. The separation follows the RATS architecture: evidence is appraised by a verifier, and the relying party applies its own policy to the resulting claims \citep{rats}.

The key broker releases only when the verified evidence and manifest policy agree; Section~\ref{sec:after} records where the evaluated version compared policy with unsigned fields instead. When the broker requires channel binding, as the reference server does, release material is sealed to the attested transport key and no plaintext model key crosses the channel. The library's default configuration does not require channel binding or cryptographic quote verification and returns the key directly. Confidential Containers Trustee and NVIDIA reference architectures already demonstrate the underlying attestation-and-secret-delivery pattern \citep{trustee,nvidia2026}. WCM standardizes the model-specific manifest, joins, lifecycle, and evidence surrounding that pattern.

\subsection{Renewable custody}

An authorization lease expires independently of process lifetime. Before expiry the runtime must obtain a new challenge and satisfy current release policy. Failure or revocation transitions the reference state machine to a terminal condition: inference is refused and the in-memory key is overwritten on a best-effort basis. The signed runtime-record chain makes the claimed sequence tamper-evident and independently replayable. It does not prove language-runtime or silicon-level zeroization.

\subsection{Derivative lineage}

Fine-tunes, adapters, merges, and quantized variants may encode economically valuable information from a protected parent. WCM records parent digests, transformation identity, custody policy, and resulting artifact digest. The lineage is evidence of declared derivation and policy continuity, not a proof that an undeclared derivative does not exist.

\section{Conformance and implementation}

The reference project includes a JSON schema frozen at version 1, a Python SDK, a reference key-broker service, hardware evidence adapters, and language-neutral conformance vectors. Ninety-one vectors are organized into four levels: document integrity and signatures; attestation-gated release scenarios; renewable runtime-custody scenarios; and lineage and accountability records. Each reportable error code is exercised. The runner states two important limitations on every execution: its portable PKI is synthetic rather than a set of live vendor roots, and it contains no GPU-side cryptographic vectors.

The implementation checks RFC 8785 canonicalization, certificate and signature chains, challenge freshness, approved measurements, transport-key binding, revocation, threshold policy, transparency receipts, derivative lineage, and terminal runtime records. Provider adapters cover SEV-SNP, TDX, Azure vTPM, and NVIDIA evidence surfaces. Integration profiles connect WCM to Confidential Containers Trustee, NVIDIA NRAS, cloud confidential-computing services, OCI, in-toto, Kyverno, vLLM, Triton, OpenTelemetry, Agent Manifest, and TRACE. These adapters are evidence of composability; they are not all equivalent hardware validations.

\section{Evaluation}

We evaluate three questions.

\textbf{E1: Does the portable verifier fail closed?} The conformance suite includes valid vectors and mutations across schema, signature, digest, freshness, policy, evidence, release, renewal, and lineage boundaries. On September 3, 2026, the complete Python suite at WCM 0.28.0 produced 618 passing tests and three skips; a rerun at the same commit on September 27, 2026 produced the same counts, and the reference conformance self-test passed all 91 vectors (L1 32, L2 37, L3 12, L4 10). The three skipped tests exercise symbolic links in artifact digests and were skipped because the Windows test account lacked the privilege to create them; they are not counted as passes. The original manuscript described them as hardware-dependent, which was incorrect. This result checks the behavior the vectors encode. It does not show that the vectors encoded every release requirement: Section~\ref{sec:after} describes requirements they omitted.

\textbf{E2: Is custody state observable after release?} Tests exercise successful release, renewal, lapse, refusal, and terminal signed records. A recorded protected-runtime experiment executed a lease to lapse inside an Azure SEV-SNP guest, refused subsequent service, wiped the Python-held key buffer, and produced a chain that verifies offline. The bounded conclusion is that the reference state machine and evidence chain behaved as specified. It does not establish physical non-recoverability.

\textbf{E3: Can the design expose rather than hide the hardware boundary?} A retained receipt dated August 27, 2026 reports a 256 MiB memory sweep inside a protected guest. Its nonce is a digest of an HCL report; it was not shown to be a fresh challenge issued by the key broker, and the raw report was not retained. Detection of controlled alias mappings is shown only in separate software tests. Neither the receipt nor those tests demonstrate resistance to an external memory interposer or an end-to-end, hardware-attested link from sweep to key release. The original manuscript described this receipt as bound to live vTPM evidence; that was not established.

\begin{table}[h]
\centering
\begin{tabular}{lll}
\toprule
Property & Evidence & Bounded conclusion \\
\midrule
Portable behavior & 91 vectors; 618 pass, 3 skip (0.28.0) & behavior the vectors encode \\
Fresh release & nonce and transport-key tests & replay/substitution rejected \\
Renewal/lapse & custody state-machine tests & terminal refusal is reproducible \\
Hardware run & signed SEV-SNP runtime records & one measured run, not certification \\
Physical owner & explicit residual-risk field & accountability, not impossibility \\
\bottomrule
\end{tabular}
\caption{Evaluation evidence and its permitted interpretation.}
\end{table}

\section{Findings after the evaluation}
\label{sec:after}

Findings published after September 3, 2026 narrow the conclusions above. They are recorded in the project's security advisory, changelog and limitations document. The first two below were reproduced for this edition by running the fixing releases' regression tests against the 0.28.0 source; those tests fail there.

\textbf{The approved workload measurement was not read from the signed report.} In the evaluated version the quote verifier received the expected workload measurement and discarded it, so the serving-image check compared the manifest with \texttt{serving\_image\_measurement}, a structured field that whoever builds the evidence controls. With channel binding and cryptographic quote verification both required, a correctly signed SEV-SNP report whose launch measurement is not approved still received the key, sealed to the presenter's own transport key, when that field named an approved value. Version 0.29.0 (September 30, 2026) compares the digest of the signed launch measurement. The same release binds a renewal decision to the evidence and challenge it answered; before it, a renewal issued to another enclave holding the same model and policy renewed this one.

\textbf{Other verification defects fixed in 0.28.5.} One key could satisfy two required signing roles, because role and signer sit outside the signed pre-image; an end-entity certificate under a trusted root could issue a further certificate that chain verification accepted; the TDX path did not check that the quoting enclave was Intel's; and two concurrent presentations of one challenge could both succeed. Version 0.28.5 (September 26, 2026) fixes each. The E1 counts above remain correct for 0.28.0, but they did not detect these defects.

\textbf{GPU confidential-compute mode was not a release condition.} Advisory GHSA-j665-99rh-w85h (high severity, published September 25, 2026) reported that in versions up to 0.28.2 the key broker populated the GPU report's confidential-compute mode but no release gate read it, and the NVIDIA provider defaulted a missing value to enabled. A confidential-computing-capable GPU with that mode disabled could therefore receive the model key. The specification had the same gap. Version 0.28.4 adds a signed requirement and a corresponding error code. The NVIDIA attestation-tool adapter still reported the mode as enabled unconditionally until 0.28.5, which makes it report the mode as unknown so that the gate denies. The evaluated version, 0.28.0, is affected, so the fail-closed result in E1 does not extend to GPU mode.

\textbf{The measured SEV-SNP platform lacks ciphertext hiding.} The cryptographic-custody claim against a hypervisor-privileged operator depends on SEV-SNP ciphertext hiding, because without it ciphertext side channels can expose keys without physical access \citep{snpciphertext}. The Azure confidential VM used for E2 and E3 reports \texttt{PLATFORM\_INFO} 0x25: alias checking complete, ciphertext hiding not enabled. A later GCP SEV-SNP capture reports the same state. On these platforms the claim holds only against adversaries without hypervisor privilege. The evaluated version did not read this field; version 0.28.2 and later can require it in signed policy.

\textbf{Attestation-key revocation is weak in practice.} A check made on September 11, 2026 found that the captured NVIDIA device chain carries certificates valid to 9999 with empty long-lived revocation lists, that AMD VCEK certificates carry serial number zero, and that an operator has no documented path to revoke a specific device key. In the evaluated version the revocation gate compared an operator-supplied set of revoked keys with a cache age the evidence reported about itself, so it could not fail unless an operator filled that set. Version 0.29.0 adds live, nonce-bound OCSP for the NVIDIA device chain that fails closed when no answer arrives; it does not establish the revocation status of the leaf the GPU issues to itself. A passing revocation check shows only that the vendor has published no revocation.

\section{Related work and novelty boundary}

Confidential computing for machine learning is a mature research area; prior systems protect models or data within TEEs, and the literature documents both security benefits and platform limitations \citep{mlccsurvey,itx}. Trustee provides reusable attestation, reference-value, KBS, and secret-delivery components \citep{trustee}. NVIDIA's architecture applies those concepts directly to proprietary models \citep{nvidia2026}. OpenSSF Model Signing authenticates model artifacts \citep{oms}; in-toto already supplies signed artifact subjects, predicates, and supply-chain provenance \citep{intoto}; and CryptoTensors applies encryption and access policy to tensor artifacts \citep{cryptotensors}. WCM neither replaces nor precedes these systems.

The narrower contribution is a vendor-neutral custody object and conformance boundary spanning exact artifact identity, workload authorization, freshness, channel-bound release, renewable post-release state, terminal evidence, and derivative lineage, with a first-class downgrade for physical ownership. This claim is falsifiable: discovery of a prior portable specification and interoperable implementation covering that complete conjunction would require the contribution to be narrowed.

\section{Limitations and conclusion}

WCM remains pre-1.0. The reference implementation is not a production certification. Vendor evidence quality varies, physical attacks can invalidate assumed roots of trust, and signed records prove only what their trusted producers could observe. The implementation's zeroization behavior is best effort at the Python layer. Independent implementations and external conformance runs are still needed.

WCM turns a deployment recipe into a portable, testable custody contract without overstating what present hardware can guarantee. Its principal result is the lifecycle join: which exact weights were authorized for which measured workload, through which fresh and channel-bound release, under which renewable policy, with what terminal and derivative evidence, and under which explicitly bounded adversary model.

\bibliographystyle{plainnat}
\bibliography{refs}
\end{document}
